% Sections 4.1 to 4.4, translated from sv/chapters/04/theory.tex.

\section{What is a system of equations?}
\label{c4:sec:ekvationssystem}
A \textbf{linear system of equations} consists of several equations in several unknowns that must
be satisfied \emph{simultaneously}. The most common is a system of \textbf{two equations in two
unknowns}:
\[
  \begin{cases}
    a_1 x + b_1 y = c_1 \\
    a_2 x + b_2 y = c_2
  \end{cases}
\]
The \textbf{solution} is the pair $(x, y)$ that satisfies both equations. Graphically, it is the
\textbf{point of intersection} of the two straight lines.

\subsection{Number of solutions}
\label{c4:sec:antal}

\begin{narrowtable}{@{}ll@{}}
  \thead{Situation} & \thead{Number of solutions} \\
  \midrule
  The lines intersect & Exactly one solution \\
  The lines are parallel (but do not coincide) & No solution \\
  The lines coincide & Infinitely many solutions \\
\end{narrowtable}

\section{The substitution method}
\label{c4:sec:substitution}
The substitution method works like this:
\begin{enumerate}
  \item Choose one equation and express one unknown in terms of the other.
  \item Substitute the expression into the other equation.
  \item Solve the resulting equation, which has one unknown.
  \item Calculate the other unknown.
  \item Check in \emph{both} of the original equations.
\end{enumerate}

\begin{exempel}
Solve the system
\[
  \begin{cases}
    x + 2y = 7 \\
    3x - y = 7
  \end{cases}
\]
\textbf{Step 1:} Equation (1) gives $x = 7 - 2y$.
\par\noindent\textbf{Step 2:} Substitute into equation (2):
\[
  3(7 - 2y) - y = 7
\]
\textbf{Step 3:} Simplify and solve:
\[
  21 - 6y - y = 7 \quad \Rightarrow \quad 21 - 7y = 7 \quad \Rightarrow \quad y = 2
\]
\textbf{Step 4:} Calculate $x$:
\[
  x = 7 - 2 \cdot 2 = 3
\]
\textbf{Step 5:} Check: $3 + 4 = 7$~\checkmark{} and $9 - 2 = 7$~\checkmark
\par\noindent The solution is $(x, y) = (3, 2)$.
\end{exempel}

\section{The elimination method (addition method)}
\label{c4:sec:addition}
\textbf{Idea:} Multiply the equations by suitable numbers so that one unknown gets coefficients
that are \emph{equal but opposite}, then add the equations to eliminate that unknown.
\begin{enumerate}
  \item Multiply one or both equations by suitable constants.
  \item Add the equations line by line, so that one unknown disappears.
  \item Solve the remaining equation.
  \item Calculate the eliminated unknown from one of the original equations.
\end{enumerate}

\begin{exempel}
Solve the system
\[
  \begin{cases}
    2x + 3y = 12 \\
    4x - y = 10
  \end{cases}
\]
\textbf{Step 1:} Multiply equation (2) by $3$:
\[
  \begin{cases}
    2x + 3y = 12 \\
    12x - 3y = 30
  \end{cases}
\]
\textbf{Step 2:} Add the equations:
\[
  14x = 42 \quad \Rightarrow \quad x = 3
\]
\textbf{Step 3:} Substitute $x = 3$ into equation (1):
\[
  6 + 3y = 12 \quad \Rightarrow \quad y = 2
\]
The solution is $(x, y) = (3, 2)$.
\end{exempel}

\section{Application: Kirchhoff's laws}
\label{c4:sec:kirchhoff}
Systems of equations are fundamental to circuit theory. With Kirchhoff's laws, the distribution of
currents and voltages in a circuit can be calculated.

\subsection{Kirchhoff's voltage law (KVL)}
\label{c4:sec:kvl}
The sum of all the voltages around a closed loop is zero:
\[
  \sum U = 0
\]

\subsection{Kirchhoff's current law (KCL)}
\label{c4:sec:kcl}
The sum of the currents into a node is zero, that is, the currents in equal the currents out:
\[
  \sum I = 0
\]

\begin{exempel}
Consider a circuit with two loops and the currents $I_1$ and $I_2$. KVL gives
\[
  \begin{cases}
    10 = 2I_1 + 4(I_1 - I_2) \\
    0 = 4(I_2 - I_1) + 6I_2
  \end{cases}
\]
which simplifies to
\[
  \begin{cases}
    6I_1 - 4I_2 = 10 \\
    -4I_1 + 10I_2 = 0
  \end{cases}
\]
Equation (2) gives $I_1 = 2.5 I_2$. Substituting this into equation (1) gives
\begin{gather*}
  6 \cdot 2.5 I_2 - 4I_2 = 10 \quad \Rightarrow \quad I_2 = \frac{10}{11} \approx 0.91\,\text{A},
  \\
  I_1 = 2.5 \cdot \frac{10}{11} = \frac{25}{11} \approx 2.27\,\text{A}.
\end{gather*}
\end{exempel}
